\section{A factorization of the implemented cobordism algebra}
\label{sec:algebra}

\subsection{Coefficient tables and gluing plans}

Let $a,b$ be loopless matchings on $2m$ exposed endpoints. Gluing $a$ to
the reflected $b$ produces $k(a,b)$ circles, where $k(a,b)\leq m$. A basis
of $\Hom(a,b)$ is indexed by the subsets of those circles: a selected circle
carries a dot. Thus a morphism is a coefficient table of length
$2^{k(a,b)}$. The program stores this table as an integer
\[
 f=\sum_{S\subseteq[k(a,b)]} f_S\,2^{\operatorname{mask}(S)},
 \qquad f_S\in\{0,1\}.
\]
The exponent in this formula is a subset mask. In particular, the bit
length of the packed integer is exponential in the number of circles.

For a composition
\[
 a\xrightarrow{f}b\xrightarrow{g}c
\]
write $k_L=k(a,b)$, $k_R=k(b,c)$, and $k_O=k(a,c)$. The existing geometry
compiler decomposes the glued surface into connected components. For a
genus-zero component $s$ it records
\[
 (L_s,R_s,B_s,e_s),
\]
where $L_s$, $R_s$, and $B_s$ are the sets of left input, right input, and
output circles belonging to the component, and $e_s$ counts any already
present dots. In ordinary composition $e_s=0$. The generalized notation is
also useful for crossing-transfer plans with cups and caps.

Within each of the three families the sets are pairwise disjoint and cover
the corresponding circle set. Empty sets are allowed. A positive-genus
component gives zero in this characteristic-two theory: the handle
operator is comultiplication followed by multiplication, and equals
multiplication by $2x=0$. Two dots on one component also give zero. These
are exactly the relations used by the baseline~\cite{barnatan,proveit}.

The old composition routine enumerates every supported pair $(S,T)$,
checks the number of dots on every component, and reconstructs its output.
Even when most pairs vanish, it has already visited them. If $a=b=c$ has
$m$ arcs, each input can have $2^m$ supported monomials, producing a
$\Theta(4^m)$ term-pair loop.

\subsection{Squarefree algebras and contraction}

For a nonnegative integer $r$ let
\[
 R_r=\F[y_1,\ldots,y_r]/(y_1^2,\ldots,y_r^2).
\]
Write $y_S=\prod_{s\in S}y_s$. Then
\begin{equation}
 y_Sy_T=
 \begin{cases}
 y_{S\cup T},&S\cap T=\varnothing,\\
 0,&S\cap T\ne\varnothing.
 \end{cases}
 \label{eq:squarefree}
\end{equation}
The dot algebra $A=\F[x]/(x^2)$ is $R_1$.

Suppose a genus-zero plan has $r$ connected components. Define
\[
 \rho_L:R_{k_L}\longrightarrow R_r,
 \qquad x_j\longmapsto y_s\quad(j\in L_s),
\]
and similarly define $\rho_R$ using the sets $R_s$. These are well-defined
algebra homomorphisms because every image variable squares to zero.

\begin{lemma}[Input contraction]
\label{lem:contraction}
If an input monomial selects two circles belonging to one component,
its image under $\rho_L$ is zero. Otherwise its image selects exactly the
components containing a chosen circle. Coefficients of input monomials
with the same image add modulo two.
\end{lemma}
\begin{proof}
Substitute the defining image of every variable into the monomial. A
repeated component introduces a factor $y_s^2$. If there is no repetition,
the product is the squarefree monomial on the selected component set.
Linearity gives the coefficient assertion.
\end{proof}

This contraction can eliminate a large fraction of a dense input before
the two morphisms are multiplied. It also combines distinct monomials by
parity, which cannot be reproduced merely by forgetting duplicate
component labels in a list of supported terms.

\subsection{The output map}

For a component with $b$ output circles, define a linear map
$T_b:A\to A^{\otimes b}$ by
\begin{align}
 T_0(1)&=0, &T_0(x)&=1,\label{eq:trace}\\
 T_b(1)&=\sum_{j=1}^{b}\prod_{\ell\ne j}z_\ell,
 &T_b(x)&=\prod_{\ell=1}^{b}z_\ell\qquad(b\geq1).
 \label{eq:comultiplication}
\end{align}
The empty product is one. Formula~\eqref{eq:comultiplication} follows by
iterating $\Delta(1)=1\otimes x+x\otimes1$ and
$\Delta(x)=x\otimes x$. The output variable sets belonging to distinct
components are disjoint.

Let $T=\bigotimes_{s=1}^{r}T_{|B_s|}$, with output variables placed at the
indices specified by the $B_s$. Multiplication by $y_s^{e_s}$ accounts for
the dots already supplied by the plan.

\begin{theorem}[Exact composition factorization]
\label{thm:factorization}
For any genus-zero compiled plan with the partition properties above,
the baseline composition is exactly
\begin{equation}
 g\circ f
 =T\left(\rho_L(f)\rho_R(g)\prod_{s=1}^{r}y_s^{e_s}\right).
 \label{eq:factorization}
\end{equation}
If the plan contains a positive-genus component, both sides are understood
to be zero. If any $e_s\geq2$, the result is zero.
\end{theorem}
\begin{proof}
Both constructions are bilinear in $f$ and $g$, so it suffices to take
monomial inputs. For each component, the total number of dots is
\[
 d_s=|S\cap L_s|+|T\cap R_s|+e_s.
\]
The contractions and multiplication in $R_r$ kill the monomial if some
$d_s\geq2$. This is the two-dot relation used by the baseline. Otherwise
the component is undotted or singly dotted. If it has no output circle,
the counit~\eqref{eq:trace} kills precisely the undotted case. If it has
output circles,~\eqref{eq:comultiplication} is precisely the baseline's
sum over leaving one output circle undotted, or its single all-dotted
term. Finally, different components use disjoint output variables, so their
contributions multiply without any additional identification. These cases
exhaust the baseline evaluation of the monomial pair. Extend by
bilinearity.
\end{proof}

\begin{center}
\begin{tikzpicture}[node distance=1.15cm,>=Latex,
 every node/.style={font=\small,align=center},
 box/.style={draw=navy!60,rounded corners,fill=blue!3,inner sep=7pt}]
\node[box] (inputs) {Two coefficient tables\\$f\in R_{k_L},\ g\in R_{k_R}$};
\node[box,below=of inputs] (contract) {Identify input dots by surface component\\$\rho_L(f),\ \rho_R(g)\in R_r$};
\node[box,below=of contract] (multiply) {Ranked subset convolution\\$H=\rho_L(f)\rho_R(g)$};
\node[box,below=of multiply] (output) {Trace or comultiply on each component\\$T(H\prod_s y_s^{e_s})\in R_{k_O}$};
\draw[->] (inputs)--(contract);
\draw[->] (contract)--(multiply);
\draw[->] (multiply)--(output);
\end{tikzpicture}
\end{center}

\subsection{Output reconstruction without exponential repeated expansion}

Expanding~\eqref{eq:comultiplication} separately for every supported
component monomial can itself repeat a large amount of work. There is a
better way: determine the one input coefficient, if any, that contributes
to a given output coefficient.

Fix an output subset $U\subseteq[k_O]$. For each component define the
required final dot count
\[
 d_s(U)=
 \begin{cases}
 1,&B_s=\varnothing,\\
 1,&B_s\subseteq U,\\
 0,&|B_s\setminus U|=1,\\
 \text{invalid},&|B_s\setminus U|\geq2.
 \end{cases}
\]
The first case is the trace condition, not a statement that a closed
component has an output dot. Subtract the fixed extra-dot count and put
$v_s(U)=d_s(U)-e_s$. If any requirement is invalid or any $v_s(U)$ lies
outside $\{0,1\}$, set the output coefficient to zero. Otherwise let
$D(U)=\{s:v_s(U)=1\}$.

\begin{lemma}[Unique coefficient lookup]
\label{lem:output}
Writing $H=\rho_L(f)\rho_R(g)$, the coefficient of $z_U$ in
\eqref{eq:factorization} is $H_{D(U)}$ when the requirements are valid, and
zero otherwise.
\end{lemma}
\begin{proof}
For a nonempty output block, a singly dotted input gives exactly the
all-dotted monomial. An undotted input gives each monomial with exactly
one missing dot, once. These alternatives have disjoint supports. For an
empty block the trace requires a singly dotted input. Thus $U$ uniquely
determines the final dot count on every component. Subtracting the dots
already supplied determines the unique contributing component monomial.
There is no multiplicity coefficient to compute.
\end{proof}

\begin{example}[Why dot identification requires parity]
Let two left circles belong to one component. The left morphism
$x_1+x_2+x_1x_2$ contracts to $y+y+0=0$. Treating its two linear terms as
one supported term would give the wrong result. Conversely, two inputs
each containing a dot on the same component multiply to zero even if their
original circle indices are different.
\end{example}

\begin{example}[Closed components and extra dots]
A closed component with no extra dot requires the coefficient of $y$;
with one extra dot it requires the coefficient of $1$. With two extra
dots it annihilates the plan. For a one-circle output, the map is the
identity on $A$ before extra dots: an undotted component outputs $1$,
not $x$.
\end{example}

\section{Ranked subset convolution and the coefficient bound}
\label{sec:convolution}

\subsection{The multiplication problem}

By~\eqref{eq:squarefree}, multiplying two coefficient arrays $F,G$ in
$R_r$ means computing
\begin{equation}
 H_S=\sum_{A\subseteq S}F_A G_{S\setminus A}.
 \label{eq:subsetconvolution}
\end{equation}
Enumerating all ordered disjoint pairs takes $3^r$ iterations. The
baseline's general-purpose pair loop can do even more work by visiting
overlapping pairs before rejecting them. Ranked subset convolution computes
all coefficients using $O((r+1)^2 2^r)$ field operations~\cite{bhkk}.
For completeness, the exact characteristic-two specialization used here
has the following short verification.

For every $S\subseteq[r]$ form rank polynomials
\[
 \widehat F_S(t)=\sum_{A\subseteq S}F_A t^{|A|},
 \qquad
 \widehat G_S(t)=\sum_{B\subseteq S}G_B t^{|B|}.
\]
Multiply these polynomials and apply subset M\"obius inversion to the
array indexed by $S$. In characteristic two the M\"obius transform is
the same XOR transform as the zeta transform. Call the resulting polynomial
array $W_S(t)$. Then
\begin{equation}
 H_S=[t^{|S|}]W_S(t).
 \label{eq:rankextraction}
\end{equation}

\begin{proof}[Verification of~\eqref{eq:rankextraction}]
Before inversion, the term $F_A G_B t^{|A|+|B|}$ occurs at every index
containing $A\cup B$. After inversion it occurs only at index
$S=A\cup B$. Requiring its degree to equal $|S|$ gives
$|A|+|B|=|A\cup B|$, which is equivalent to disjointness. Thus the
remaining terms are exactly those of~\eqref{eq:subsetconvolution}.
\end{proof}

Only ranks $0,\ldots,r$ are needed. Higher terms may be discarded before
the inverse transform because that transform does not mix polynomial
degrees. An unranked zeta transform would compute union convolution and
would incorrectly replace $x\cdot x=0$ by $x\cdot x=x$.

\subsection{Why the number of surface components stays small}

\begin{lemma}[Component count for matching composition]
\label{lem:components}
For a composition of three matchings on $2m$ endpoints,
\[
 k_L,k_R,k_O\leq m,
 \qquad r\leq\min(k_L,k_R)\leq m.
\]
\end{lemma}
\begin{proof}
Every circle in a pair of matchings uses at least one arc from either
matching, giving the first three bounds. Every left input disc meets at
least one middle-matching arc, and so is glued to a right input disc. The
same argument applies on the right. Consequently each connected glued
component contains at least one left disc and one right disc. Distinct
components contain disjoint input discs, giving the asserted bound on $r$.
\end{proof}

This argument is needed. For an arbitrary abstract list of surface
components with many empty input blocks, $r$ can exceed the three circle
counts. The generic implementation supports such plans but charges for
$2^r$ explicitly; the matching-composition bound is applied only after
Lemma~\ref{lem:components}.

\begin{theorem}[Dense local coefficient complexity]
\label{thm:densecost}
Given a compiled composition plan of matchings on $2m$ endpoints and the
two full coefficient tables over $\F$, the output can be computed in
$\poly(m)2^m$ bit operations and $\poly(m)2^m$ bits of working storage.
The field-operation count of the ranked multiplication is
$O((m+1)^2 2^m)$. The implementation's explicit array and index arithmetic
admits a conservative bound $O((m+1)^3 2^m)$ for the coefficient work in a
random-access bit-cost model. Label manipulation in constructing the plan
adds a polynomial in $m$ and the label bit length.
\end{theorem}
\begin{proof}
For each input, precompute the image of every subset by taking off one
selected variable at a time. Store a sentinel for a repeated component.
This uses $O(2^{k_L}+2^{k_R})$ entries with $O(m)$ bits per entry.
Read each coefficient once, and toggle its contracted coefficient if its
image is valid. Ranked convolution uses arrays of $2^r$ polynomials, each
of at most $r+1$ bits. Its zeta transforms perform $O(r2^r)$ XORs on
$O(r)$-bit words. Naive carryless multiplication of two rank polynomials
takes $O((r+1)^2)$ bit operations, at each of $2^r$ indices.

For every one of $2^{k_O}$ output subsets, inspect its intersection with
each output block and use Lemma~\ref{lem:output}. This uses only integers
of $O(m)$ bits. The input and output full morphisms are converted to and
from byte arrays once. In particular, there is no loop making $2^m$
updates to an integer of $2^m$ bits. Charging conservatively for array
indices and small masks adds at most an additional polynomial factor.
Lemma~\ref{lem:components} bounds all four exponential dimensions by
$2^m$. The stored tables have $O(2^m)$ entries of $O(m)$ bits, together
with the packed inputs and output.
\end{proof}

\begin{remark}[Meaning of the improvement]
The theorem compares complete coefficient operations. The old
$\Theta(4^m)$ statement counts term pairs; it is not an assertion that
the old bit complexity is exactly $\Theta(4^m)$. Its packed-integer
operations can cost more. The new theorem supplies a bit bound for its
own coefficient kernel. On a model that explicitly writes the output
table, the worst-case output size already costs $\Omega(2^m)$, so the
remaining overhead is polynomial in $m$.
\end{remark}

\begin{example}[An explicit dense test family]
In $\End(a)=R_m$, take
\[
 F=\prod_{i=1}^m(1+x_i),\qquad G=F+1.
\]
The two supports have sizes $2^m$ and $2^m-1$. Since every positive-degree
element squares to zero in this commutative characteristic-two algebra,
$F^2=1$ and $FG=F+1$. The baseline nevertheless visits
$2^m(2^m-1)$ pairs. This is a family of valid algebra inputs, not a claim
that the scanner must produce these particular morphisms from knot
diagrams.
\end{example}
